% !PRINTED 100
\weberpage{100}
%<<CONT
der Zahlen 0, 1, 2 \dots $a - 1$ enthalten ist, und dadurch ist dann
die Substitution $S_0$ \gesperrt{vollständig bestimmt}.

Wenn die vier Transformationszahlen einen gemeinsamen
Faktor haben, so läßt sich dieser mittels der Formel
\[
\begin{pmatrix}m, & 0\\0, & m\end{pmatrix}\begin{pmatrix}a, & b\\c, & \partial\end{pmatrix} = \begin{pmatrix}ma, & mb\\mc, & m\partial\end{pmatrix}
\]
durch Zusammensetzung mit der Multiplikation absondern, und
wir setzen demnach jetzt voraus, daß $a, b, c, \partial$ \gesperrt{keinen gemeinschaftlichen
 Teiler haben}. Man kann dann immer die zwei
ganzen Zahlen $\xi, \eta$ so bestimmen, daß
\weq{3}{\begin{gathered} a\eta - c\xi = \alpha\\ b\eta - \partial\xi = \beta \end{gathered}}
ohne gemeinsamen Teiler sind.

Um dies einzusehen, setzen wir zunächst $\xi, \eta$ relativ prim
voraus. Dann ist jeder gemeinsame Teiler von $\alpha, \beta$ notwendig
Teiler von $n$, wie man aus den Auflösungen von (3)
\weq{4}{\begin{gathered} n\xi = b\alpha - a\beta\\ n\eta = \partial\alpha - c\beta \end{gathered}}
erkennt. Nimmt man also $\xi$ \gesperrt{nicht teilbar} durch alle in $a$
und $b$ zugleich aufgehenden Primzahlen, dagegen $\xi$ \gesperrt{teilbar},
$\eta$ \gesperrt{unteilbar} durch alle anderen in $n$ aufgehenden Primzahlen,
und überdies $\xi, \eta$ relativ prim, was stets möglich ist, so haben
$\alpha$ und $\beta$ keinen gemeinsamen Teiler. Hierauf bestimmt man $\gamma, \delta$
so, daß
\[
\alpha\delta - \beta\gamma = 1.
\]

Es ist dann nach (3) auch
\[
(a\delta - b\gamma)\eta - (c\delta - \partial\gamma)\xi = 1,
\]
und es ergibt sich die folgende Zusammensetzung, wie leicht mit
Benutzung von (4) erkannt wird:
\weq{5}{\begin{pmatrix}a, & b\\c, & \partial\end{pmatrix} = \begin{pmatrix}a\delta - b\gamma, & \xi\\c\delta - \partial\gamma, & \eta\end{pmatrix}\begin{pmatrix}1, & 0\\0, & n\end{pmatrix}\begin{pmatrix}\alpha, & \beta\\\gamma, & \delta\end{pmatrix}.}

Nennen wir also
\weq{6}{\begin{pmatrix}1, & 0\\0, & n\end{pmatrix}}
die \gesperrt{Haupttransformation} vom Grade $n$, so ist damit bewiesen,
\gesperrt{daß sich alle Transformationen vom Grade $n$ aus einer
Multiplikation, einer Haupttransformation und linearen
Transformationen zusammensetzen lassen.}
